% =====================================================================
% Capítulo: Diseño del Sistema
% Fuentes: borradores/05-diseno.tex, docs/tfg/DECISIONES.md (ADR-001…021),
% AGENTS.md, apps/web/AGENTS.md, apps/web/supabase/AGENTS.md,
% apps/web/src/config/*.config.ts, packages/cms/*,
% packages/function-middleware. Estado de las fuentes: F3c cerrada.
% =====================================================================

Este capítulo describe el diseño de PymeKit: la arquitectura general, el
modelo multi-tenant y de control de acceso, la estrategia de seguridad en
profundidad, la integración con la pasarela de pago, el diseño de la interfaz
y los mecanismos de configuración que hacen la plataforma reutilizable. Cada
decisión remite al registro de decisiones de arquitectura (ADR,
\textit{Architecture Decision Record}~\cite{nygard2011adr}) del proyecto, que
se resume en la sección~\ref{sec:dis-decisiones}. La infraestructura de
ejecución (contenedores, nube y puesta en marcha) se trata en el capítulo de
Despliegue del Sistema.

% ---------------------------------------------------------------------
\section{Arquitectura general}
\label{sec:dis-arquitectura}

\subsection{Monorepo de paquetes}
\label{sec:dis-monorepo}

El código se organiza en un único \textit{monorepo} (un repositorio con
varias aplicaciones y paquetes que se versionan juntos) gestionado con
Turborepo~\cite{turborepo} y pnpm (ADR-002):

\begin{itemize}
  \item \texttt{apps/web}: la aplicación web, que incluye la web pública, el
        área privada, la consola de administración y el CMS;
  \item \texttt{apps/web/supabase}: esquemas, migraciones, datos de prueba y
        pruebas pgTAP de \emph{toda} la base de datos;
  \item \texttt{apps/e2e}: las pruebas E2E con Playwright~\cite{playwright};
  \item \texttt{packages/*}: los paquetes reutilizables de la plataforma
        (\texttt{@pymekit/*}: interfaz, cliente de base de datos,
        \textit{middleware}, autenticación, cuentas, equipos, facturación,
        administración, i18n\ldots);
  \item \texttt{packages/cms/*}: los paquetes del CMS
        (\texttt{@pymekit/cms-*}), separados en paquetes de servidor
        (servicios y rutas de la API) y paquetes de cliente (componentes de
        las pantallas).
\end{itemize}

El monorepo contiene 51 paquetes, 27 de ellos del CMS, sin ciclos de
dependencias (RNF-04). Cada paquete declara sus dependencias de forma
explícita y separa las exportaciones de cliente y de servidor. Esta
separación evita que código con credenciales de servidor acabe en el
navegador ---la \textit{build} de la web falla si el cliente de base de datos
del servidor aparece en un fragmento destinado al navegador--- y, como se
comprobó en la incidencia B-16, también evita ciclos entre paquetes.

\subsection{Capas y pila tecnológica}
\label{sec:dis-capas}

La aplicación sigue una arquitectura en capas en la que cada capa solo
depende de la inferior y la base de datos es la última barrera de
autorización:

\begin{enumerate}
  \item \textbf{Presentación.} Rutas por fichero de TanStack
        Router~\cite{tanstack-router} con React 19, renderizadas en el
        servidor (SSR). Componentes accesibles de \texttt{@pymekit/ui}
        con Tailwind CSS. La caché de datos del cliente usa TanStack
        Query~\cite{tanstack-query}.
  \item \textbf{Aplicación.} \textit{Server functions} de TanStack
        Start~\cite{tanstack-start}: funciones que se ejecutan en el servidor
        y se invocan desde el cliente con tipos compartidos, validadas con
        Zod y protegidas por una cadena de \textit{middleware} (RF-02). Para el
        CMS, una API HTTP construida con Hono~\cite{hono}.
  \item \textbf{Servicios.} Lógica de negocio en servicios sin dependencia de
        la interfaz, con el cliente de datos inyectado: el cliente de
        Supabase~\cite{supabase-docs} en la web y Drizzle ORM~\cite{drizzle}
        en el CMS.
  \item \textbf{Datos.} PostgreSQL con políticas de seguridad a nivel de fila
        (RLS), funciones y \textit{triggers}; servicio de autenticación con
        JWT y MFA; almacenamiento de ficheros.
\end{enumerate}

Para que el conjunto sea coherente se usa una sola pila de librerías en toda
la web, CMS incluido: un único sistema de formularios, uno de
internacionalización y una sola biblioteca de componentes (ADR-013). Así cada
cosa se hace de una única forma (RNF-05).

\subsection{Una única base de datos}
\label{sec:dis-bd-unica}

La aplicación y el CMS comparten \textbf{una sola instancia} de PostgreSQL
(ADR-002). Los datos de la plataforma viven en el esquema \texttt{public} y
los del CMS en un esquema propio, \texttt{cms} (ADR-012). Compartir la base de
datos es lo que permite que el CMS administre los datos de la propia
plataforma ---incluido el blog que publica la web (ADR-017)--- y que las
mismas políticas de seguridad protejan ambas partes. Los esquemas
declarativos se numeran (del \texttt{00} al \texttt{55}) para fijar el orden
de creación, y las migraciones se generan a partir de ellos.

\subsection{CMS integrado en la consola de administración}
\label{sec:dis-cms-integrado}

La decisión de mayor calado del diseño es cómo integrar el CMS (ADR-011). Un
CMS de este tipo se construye habitualmente como una aplicación de página
única con su propio enrutador y una API independiente. Se evaluaron tres opciones:

\begin{enumerate}[label=(\alph*)]
  \item \textbf{aplicaciones separadas} (la web, la interfaz del CMS y su
        API como tres servicios): la de menor esfuerzo, pero con tres
        despliegues, dos inicios de sesión y las \textit{cookies} de sesión
        duplicadas;
  \item \textbf{integración híbrida} (la API dentro de la web y la interfaz
        del CMS servida como aplicación estática): convivirían dos
        enrutadores y dos formas de programar;
  \item \textbf{integración total}: la API se monta dentro del servidor de
        la web y las pantallas se reescriben como rutas de la web.
\end{enumerate}

Se eligió la opción (c), aun siendo la de mayor esfuerzo estimado, porque da
un único servicio, un único inicio de sesión y un único estilo de código:

\begin{itemize}
  \item La aplicación Hono del CMS se monta en \texttt{/api/cms/*} mediante
        una ruta de servidor de la web que le reenvía la petición tal cual.
        El CMS comparte así proceso, dominio y \textit{cookies} con la web.
  \item Las pantallas del CMS son rutas de la web bajo \texttt{/admin/cms/*},
        con \textit{loaders} y TanStack Query sobre un cliente RPC tipado de
        la API.
  \item Durante el SSR, los \textit{loaders} no hacen una petición de red a la
        propia API: invocan la aplicación Hono dentro del mismo proceso,
        reenviando solo la cabecera \texttt{cookie} (ADR-016).
\end{itemize}

La figura~\ref{fig:dis-arquitectura} resume la arquitectura resultante.

\begin{figure}[htbp]
  \centering
  \resizebox{\textwidth}{!}{%
  \begin{tikzpicture}[
      font=\footnotesize,
      box/.style={draw, rounded corners, align=center, minimum height=0.9cm,
                  inner sep=4pt, fill=white},
      group/.style={draw, dashed, rounded corners, inner sep=8pt},
      arr/.style={-{Stealth[length=2mm]}}
    ]
    \node[box, minimum width=11cm] (browser) at (0, 5.2)
      {Navegador (React 19): web pública, área privada, consola \texttt{/admin} y CMS \texttt{/admin/cms}};

    \node[box, minimum width=3.0cm] (ssr) at (-3.8, 3.3)
      {Rutas y \textit{loaders}\\(SSR, TanStack Router)};
    \node[box, minimum width=3.0cm] (sfn) at (0, 3.3)
      {\textit{Server functions}\\+ \textit{middleware}};
    \node[box, minimum width=3.0cm] (hono) at (3.8, 3.3)
      {Ruta \texttt{/api/cms/*}\\$\rightarrow$ aplicación Hono};
    \node[box, minimum width=3.0cm] (svc) at (-1.9, 1.6)
      {Servicios + cliente Supabase\\(sesión del usuario)};
    \node[box, minimum width=3.0cm] (drizzle) at (3.8, 1.6)
      {Servicios Drizzle\\(\textit{claims} del usuario)};
    \node[box, minimum width=2.4cm] (wh) at (-6.6, 1.6)
      {Ruta \texttt{/api/billing}\\\textit{webhook}};

    \begin{scope}[on background layer]
      \node[group, fit=(ssr)(sfn)(hono)(svc)(drizzle)(wh),
            label={[font=\small\bfseries, anchor=north east]south east:\texttt{apps/web} (TanStack Start, un único proceso)}] (web) {};
    \end{scope}

    \node[box, minimum width=2.6cm] (auth) at (-3.8, -0.6) {Auth\\(JWT, MFA)};
    \node[box, minimum width=4.6cm] (pg) at (0.4, -0.6)
      {PostgreSQL + RLS\\esquemas \texttt{public}, \texttt{cms} y \texttt{auth}};
    \node[box, minimum width=2.2cm] (st) at (4.4, -0.6) {Storage};
    \begin{scope}[on background layer]
      \node[group, fit=(auth)(pg)(st),
            label={[font=\small\bfseries]below:Supabase (una única instancia)}] (sb) {};
    \end{scope}

    \node[box] (stripe) at (-10.6, 1.6) {Pasarela\\de pago};

    \draw[arr] (browser) -- (ssr);
    \draw[arr] (browser) -- (sfn);
    \draw[arr] (browser) -- node[right, font=\scriptsize]{RPC} (hono);
    \draw[arr] (ssr) -- (sfn);
    \draw[arr, dashed] (ssr.south) -- ++(0,-0.3) coordinate (ssrq) -- (ssrq -| hono.south) node[midway, below, font=\scriptsize, fill=white, inner sep=1pt]{SSR en proceso} -- (hono.south);
    \draw[arr] (sfn) -- (svc);
    \draw[arr] (hono) -- (drizzle);
    \draw[arr] (svc) -- (pg);
    \draw[arr] (svc) -- (auth);
    \draw[arr] (drizzle) -- (pg);
    \draw[arr] (stripe) -- node[above, font=\scriptsize]{\textit{webhooks}} (wh);
    \draw[arr] (wh) |- (pg.west);
    \draw[arr, dashed] (sfn.north west) to[bend right=25] node[above, font=\scriptsize, pos=0.6]{\textit{checkout}} (stripe.north);
  \end{tikzpicture}}
  \caption{Arquitectura general de PymeKit: una única aplicación web con el
    CMS integrado y una única instancia de la base de datos.}
  \label{fig:dis-arquitectura}
  % Texto alternativo: el navegador se comunica con la aplicación web, que
  % contiene las rutas con SSR, las server functions y la ruta /api/cms que
  % monta la aplicación Hono del CMS. Las server functions usan servicios con
  % el cliente de Supabase y la sesión del usuario; la API del CMS usa
  % servicios Drizzle con los claims del usuario. Ambos acceden a una única
  % instancia (Auth, PostgreSQL con RLS y Storage). La pasarela de pago
  % recibe las sesiones de pago desde las server functions y notifica los
  % cambios a la ruta de webhook, que actualiza la base de datos.
\end{figure}

% ---------------------------------------------------------------------
\section{Modelo multi-tenant y control de acceso}
\label{sec:dis-multitenant}

\subsection{Cuentas como unidad de aislamiento}
\label{sec:dis-cuentas}

PymeKit sigue un modelo multi-tenant de \emph{base de datos y esquema
compartidos}: todos los clientes comparten tablas y cada fila lleva el
identificador de la cuenta a la que pertenece. Es el modelo de menor coste de
operación y el que más simplifica las actualizaciones, a cambio de que el
aislamiento dependa por completo de la corrección del control de
acceso~\cite{bezemer2010multitenant}. Por eso el diseño traslada ese control a la base de
datos.

El \textit{tenant} es la cuenta (sección~\ref{sec:ana-modelo-datos}). La
relación de un usuario con una cuenta se resuelve con un pequeño conjunto de
funciones auxiliares sobre las que se escriben todas las políticas:
\texttt{has\_role\_on\_account} (pertenencia), \texttt{has\_permission}
(permiso concreto) y \texttt{has\_more\_elevated\_role} (jerarquía). Las
escrituras destructivas se condicionan al permiso, no a la mera pertenencia,
y la cuenta activa de cada sesión se guarda en el servidor
(\texttt{user\_active\_account}) para que el área privada sepa en qué
contexto trabaja el usuario.

El modelo admite tres modos, elegidos por configuración
(sección~\ref{sec:dis-configuracion}): solo cuentas personales (B2C), solo
cuentas de equipo (B2B) o ambas. En la web de una sola pyme (ADR-020), los
equipos pueden representar, por ejemplo, las tiendas de la empresa, y sus
clientes usar la cuenta personal.

\subsection{Control de acceso del CMS}
\label{sec:dis-rbac-cms}

El CMS tiene su propio control de acceso basado en roles
(\textit{Role-Based Access Control}, RBAC) en el esquema \texttt{cms}. Conviven
así dos modelos de autorización en la misma base de datos
(tabla~\ref{tab:dis-dos-modelos}). La decisión de diseño (ADR-014) fue
conectarlos en lugar de sustituir uno por otro:

\begin{itemize}
  \item \textbf{El super-administrador es la raíz del CMS.} Unos
        \textit{triggers} sobre los usuarios sincronizan a los
        super-administradores con un rol de sistema \emph{Root}, de rango
        máximo, que concede todos los permisos. Al promocionar a un usuario
        recibe acceso, una cuenta activa del CMS y ese rol; al retirarle la
        condición pierde el acceso y su cuenta se desactiva, pero se conserva
        para la auditoría. Un índice único garantiza que solo exista un rol
        raíz.
  \item \textbf{El RBAC del CMS da acceso limitado a otro personal}
        (soporte, gestión de contenidos). Se descartó dejar solo al
        super-administrador porque incumpliría RF-09 («con permisos
        propios»). Para que una instalación nueva no empiece con el RBAC
        vacío, una migración de datos idempotente deja preparados dos roles
        de partida, editables y sin nadie asignado (ADR-026): «Soporte»,
        que solo lee cuentas, miembros, usuarios y la auditoría, y «Gestor
        de contenidos», que gestiona el blog y los archivos a través de un
        grupo de permisos «Contenidos». Se eligió una migración y no el
        \textit{seed} porque este solo se aplica en desarrollo.
\end{itemize}

\begin{table}[htbp]
  \centering
  \footnotesize
  \caption{Los dos modelos de autorización de PymeKit.}
  \label{tab:dis-dos-modelos}
  \begin{tabularx}{\textwidth}{@{}>{\raggedright\arraybackslash}p{2.6cm}XX@{}}
    \toprule
     & \textbf{Plataforma (\texttt{public})} & \textbf{CMS (\texttt{cms})} \\
    \midrule
    Sujetos & Usuarios de las cuentas (clientes) & Personal con cuenta del CMS \\
    Ámbito & Una cuenta (\textit{tenant}) & Toda la base de datos administrada \\
    Roles & Jerarquía por \texttt{hierarchy\_level} & Rango único 0--100, con vigencia \\
    Permisos & Cinco permisos por funcionalidad & De datos (esquema, tabla, columna), de sistema y de almacenamiento; grupos; concesiones y denegaciones por cuenta \\
    Raíz & Propietario principal de cada cuenta & Rol Root, ligado al super-administrador \\
    MFA & Exigido si el usuario lo tiene configurado & Obligatorio por defecto \\
    \bottomrule
  \end{tabularx}
\end{table}

\subsection{Acceso a la consola}
\label{sec:dis-acceso-consola}

Super-administradores y personal del CMS comparten la consola, con límites
distintos (ADR-016): la consola admite a quien es super-administrador
\emph{o} tiene acceso al CMS; las páginas de gestión de la plataforma exigen
super-administrador y redirigen al personal al CMS. La barra lateral solo
muestra las secciones que la API permite, pero esa ocultación es solo de
interfaz: cada ruta y cada \textit{endpoint} vuelven a comprobar el permiso.

% ---------------------------------------------------------------------
\section{Seguridad en profundidad}
\label{sec:dis-seguridad}

El principio que guía el diseño es que \textbf{la base de datos es la
frontera de autorización} (RNF-02): toda comprobación en la interfaz o en el
servidor de aplicación es una primera barrera, útil para la experiencia de
uso, pero nunca suficiente. Si una capa superior falla, las políticas de la
base de datos siguen impidiendo el acceso, en línea con las recomendaciones
de verificación de seguridad de OWASP~\cite{owasp-asvs}.

\subsection{RLS como frontera}
\label{sec:dis-rls}

PostgreSQL permite asociar a cada tabla políticas de seguridad a nivel de
fila (\textit{Row-Level Security}, RLS)~\cite{postgresql-rls}: predicados que
el propio motor añade a cada consulta según el usuario que la ejecuta.
Supabase~\cite{supabase-rls} traslada la identidad del usuario (su JWT) a la
sesión de PostgreSQL, de modo que las políticas deciden con ella. Las reglas
de diseño, recogidas en las directrices de la base de datos, son:

\begin{itemize}
  \item RLS activo en toda tabla. Se revocan los privilegios por defecto a
        los tres roles de la API (\texttt{anon}, \texttt{authenticated} y
        \texttt{service\_role}), incluidos los que ignoran RLS como
        \texttt{TRUNCATE}, y se conceden solo los verbos que la funcionalidad
        usa.
  \item El \texttt{UPDATE} se concede por columnas, dejando fuera las de
        identidad y de pertenencia (\texttt{id}, \texttt{account\_id}), para
        que nadie pueda mover una fila a otra cuenta.
  \item Las políticas de \texttt{UPDATE} repiten en \texttt{WITH CHECK} las
        comprobaciones del \texttt{INSERT}: sin esa cláusula PostgreSQL solo
        valida la fila antigua~\cite{postgresql-createpolicy} (B-05, B-50).
  \item Toda función \texttt{security definer}, que se ejecuta con
        privilegios elevados y se salta RLS, aplica sus propias
        comprobaciones y las justifica en un comentario; las funciones
        internas no son ejecutables por los roles de la API (B-45).
  \item Las llamadas a funciones dentro de un predicado se envuelven en
        \texttt{(select \ldots)} para que se evalúen una sola vez (B-55).
  \item El cliente de servidor que ignora RLS se reserva para casos
        excepcionales, tras validar a mano los permisos.
\end{itemize}

\subsection{Segundo factor obligatorio en la administración}
\label{sec:dis-mfa}

El servicio de autenticación distingue el nivel de garantía de la sesión
(\textit{Authenticator Assurance Level}, AAL)~\cite{supabase-mfa}:
\texttt{aal1} tras la contraseña y \texttt{aal2} tras completar el segundo
factor. El diseño exige \texttt{aal2} en todas las capas que dan acceso a la
administración:

\begin{itemize}
  \item en la base de datos, la función que identifica al super-administrador
        devuelve falso si la sesión no es \texttt{aal2}; unas políticas
        restrictivas exigen \texttt{aal2} a quien tiene un factor
        configurado; y el acceso al CMS lo decide
        \texttt{cms.verify\_admin\_access()}, que comprueba la cuenta activa
        del CMS y el MFA;
  \item en la interfaz, la consola exige siempre \texttt{aal2}: con un factor
        configurado redirige a la verificación y vuelve a la página pedida;
        sin factor responde 404 (ADR-016).
\end{itemize}

\subsection{Middleware más comprobación en la base de datos}
\label{sec:dis-middleware}

La figura~\ref{fig:dis-flujo-cms} muestra el recorrido de una petición a la
API del CMS y las comprobaciones que encuentra en cada capa. Las
\textit{server functions} de la web siguen el mismo patrón: una cadena de
\textit{middleware} tipado (autenticación, pertenencia al equipo, permiso por
funcionalidad, super-administrador) valida la petición y después el cliente de
Supabase con la sesión del usuario deja que RLS tome la decisión final.

\begin{figure}[htbp]
  \centering
  \resizebox{\textwidth}{!}{%
  \begin{tikzpicture}[
      font=\footnotesize,
      paso/.style={draw, rounded corners, align=center, minimum width=3.3cm,
                   minimum height=1.0cm, inner sep=3pt},
      check/.style={align=left, font=\scriptsize, text width=7.6cm},
      arr/.style={-{Stealth[length=2mm]}}
    ]
    \node[paso] (s1) at (0, 0)     {Navegador\\(consola \texttt{/admin/cms})};
    \node[paso] (s2) at (0, -1.6)  {Ruta de servidor\\\texttt{/api/cms/\$}};
    \node[paso] (s3) at (0, -3.2)  {Aplicación Hono\\(\textit{middleware})};
    \node[paso] (s4) at (0, -4.8)  {Servicio Drizzle\\(transacción)};
    \node[paso] (s5) at (0, -6.4)  {PostgreSQL\\(esquema \texttt{cms})};

    \draw[arr] (s1) -- (s2);
    \draw[arr] (s2) -- (s3);
    \draw[arr] (s3) -- (s4);
    \draw[arr] (s4) -- (s5);

    \node[check, right=0.6cm of s1] {Guarda de ruta: sesión viva y
      \texttt{aal2}, acceso de administración (si no: inicio de sesión,
      verificación MFA o 404)};
    \node[check, right=0.6cm of s2] {Reenvío a Hono; el código del CMS se
      importa solo en el servidor};
    \node[check, right=0.6cm of s3] {Cabeceras de seguridad y CSRF propia; JWT
      válido (401), acceso al CMS (403), usuario no bloqueado (en cerrado);
      validación estricta de la entrada con Zod (400)};
    \node[check, right=0.6cm of s4] {\textit{Claims} del usuario y rol
      \texttt{authenticated} fijados con parámetros enlazados; permiso
      concreto comprobado por el servicio};
    \node[check, right=0.6cm of s5] {\texttt{cms.verify\_admin\_access()}
      (cuenta activa, MFA) y políticas RLS: \textbf{decisión final}};
  \end{tikzpicture}}
  \caption{Flujo de una petición a la API del CMS y comprobaciones de
    seguridad en cada capa.}
  \label{fig:dis-flujo-cms}
  % Texto alternativo: diagrama vertical de cinco pasos (navegador, ruta de
  % servidor, aplicación Hono, servicio Drizzle, PostgreSQL). A la derecha de
  % cada paso se indican sus comprobaciones; la última, la base de datos, es
  % la que toma la decisión.
\end{figure}

Dos detalles del diseño merecen mención. Primero, el servicio fija la
identidad del usuario en la transacción mediante parámetros enlazados, nunca
interpolando valores del JWT en el SQL, de modo que las políticas RLS del
esquema \texttt{cms} se aplican igual que a cualquier otro cliente. Segundo,
un super-administrador con sesión \texttt{aal1} supera el \textit{middleware}
(tiene el \textit{claim} de acceso), pero no ve datos, porque la base de
datos exige el segundo factor: es la defensa en profundidad en acción.

\subsection{Fallo en cerrado y errores sin detalles internos}
\label{sec:dis-fallo-cerrado}

Todas las comprobaciones de seguridad \textbf{fallan en cerrado}
(\textit{fail closed}): ante un dato ausente, ilegible o no válido, se
deniega. El diseño lo aplica a la opción que indica si el CMS exige MFA, que
se lee con una función de privilegios elevados para que RLS no pueda
ocultársela a quien decide (B-06); a la comprobación de usuario bloqueado
cuando el servicio de autenticación no responde (B-23); y a la auditoría,
cuyo fallo deshace la escritura auditada (B-41).

Del mismo modo, los errores de la base de datos se clasifican en el servidor:
el cliente solo recibe un estado HTTP y un código estable
(\texttt{RECORD\_*}, \texttt{SETTINGS\_*}, \texttt{AUDIT\_LOG\_*}\ldots),
nunca el texto de PostgreSQL, que podría revelar nombres de tablas o
restricciones (B-20, B-33).

\subsection{Seguridad del entorno de desarrollo}
\label{sec:dis-seguridad-entorno}

La seguridad alcanza también al entorno de trabajo. Los servicios locales de
la base de datos se aíslan de Internet con una regla de cortafuegos específica
para contenedores y la aplicación se consulta mediante un túnel SSH (ADR-010,
tras la incidencia B-04). Los secretos nunca se versionan.

% ---------------------------------------------------------------------
\section{Integración con la pasarela de pago}
\label{sec:dis-pagos}

La facturación se diseña detrás de una \emph{pasarela abstracta}
(\texttt{@pymekit/billing-gateway}) con una estrategia por proveedor: la
aplicación pide «crear una sesión de pago» o «abrir el portal de cliente» sin
conocer el proveedor concreto, que se elige con la variable
\texttt{VITE\_BILLING\_PROVIDER} y se valida al construir la aplicación.
Actualmente se registra Stripe.

Los planes, sus precios y sus características se declaran en
\texttt{billing.config.ts}, validado con un esquema. El flujo tiene dos
sentidos (figura~\ref{fig:dis-arquitectura}): las \textit{server functions}
crean la sesión de pago tras comprobar el permiso \texttt{billing.manage}; y
la pasarela notifica cada cambio a la ruta \texttt{/api/billing/webhook}, que
verifica la firma y sincroniza clientes, suscripciones y pedidos con el
cliente de servicio. Ningún usuario escribe esas tablas con su sesión
(RN-07): la fuente de verdad de un cobro es la pasarela.

Los planes de partida están pensados para una pyme (ADR-025): «Básico»,
«Profesional» y «Empresa», en euros, cada uno con pago mensual y anual. Los
importes son ficticios y los identificadores de precio de la pasarela se
leen de variables de entorno (\texttt{VITE\_STRIPE\_PRICE\_*}); si faltan,
se usan valores de relleno marcados como ejemplo, de modo que la página de
precios funciona desde el primer arranque y la pyme solo tiene que crear sus
precios en Stripe y copiar sus identificadores. Los textos de cada plan son
claves de traducción.
\todo{Completar con la verificación con la pasarela en modo de pruebas (F6/F7).}

% ---------------------------------------------------------------------
\section{Diseño de la interfaz}
\label{sec:dis-interfaz}

\subsection{Web pública y área privada}

La web pública ofrece inicio, precios, blog, preguntas frecuentes, contacto y
textos legales; el área privada, el panel de la cuenta activa, sus ajustes,
miembros y facturación. La identidad visual es propia (logotipo, iconos y
capturas generadas a partir de la propia aplicación, B-46). La interfaz es
adaptable a dispositivos móviles, admite tema claro y oscuro y se construye
con componentes accesibles (RNF-08).

Mientras PymeKit se presenta como TFG, la página de inicio incluye un bloque
académico (logotipo de la escuela, autor, tutor y un recorrido guiado por la
demostración). Depende de un único interruptor de configuración, de modo que
una pyme que reutilice la plataforma lo desactiva sin tocar el código
(sección~\ref{sec:dis-configuracion}).

\subsection{Consola de administración por áreas de negocio}
\label{sec:dis-consola}

La primera versión de la consola separaba «Plataforma», «CMS» y «Recursos» y
agrupaba las tablas por esquema SQL, una organización técnica que no entiende
quien gestiona una pyme. Tras la revisión del autor se rediseñó como una
\textbf{única consola organizada por áreas de negocio} (ADR-020,
figura~\ref{fig:dis-consola}):

\begin{itemize}
  \item un bloque \textbf{Gestión} con Inicio, Gestión de cuentas (solo
        super-administrador) y las herramientas del CMS: Usuarios, Archivos,
        Paneles, Auditoría y Ajustes;
  \item un bloque \textbf{Datos} con una entrada plegable por área (Blog,
        Cuentas, Facturación, Sistema y Otros datos) que contiene sus tablas
        con nombre en español, y una entrada «Todas las tablas».
\end{itemize}

El área de cada tabla, su nombre visible y su descripción son \emph{datos}
de los metadatos del CMS, editables en Ajustes $\rightarrow$ Recursos, de
modo que una pyme puede crear áreas para sus propias tablas sin programar.
Las tablas de esquemas protegidos no aparecen en la navegación.

\begin{figure}[htbp]
  \centering
  \begin{tikzpicture}[
      font=\scriptsize,
      item/.style={anchor=west, inner sep=1pt},
      head/.style={anchor=west, font=\scriptsize\bfseries, inner sep=1pt}
    ]
    \draw[rounded corners] (0,0) rectangle (4.2,-8.6);
    \node[head] at (0.2,-0.35) {Gestión};
    \node[item] at (0.35,-0.8)  {Inicio};
    \node[item] at (0.35,-1.2)  {Gestión de cuentas$^{*}$};
    \node[item] at (0.35,-1.6)  {Usuarios};
    \node[item] at (0.35,-2.0)  {Archivos};
    \node[item] at (0.35,-2.4)  {Paneles};
    \node[item] at (0.35,-2.8)  {Auditoría};
    \node[item] at (0.35,-3.2)  {Ajustes};
    \node[head] at (0.2,-3.8) {Datos};
    \node[item] at (0.35,-4.25) {$-$ Blog};
    \node[item] at (0.75,-4.65) {Entradas};
    \node[item] at (0.75,-5.0)  {Categorías};
    \node[item] at (0.75,-5.35) {Etiquetas};
    \node[item] at (0.35,-5.8)  {$+$ Cuentas};
    \node[item] at (0.35,-6.2)  {$+$ Facturación};
    \node[item] at (0.35,-6.6)  {$+$ Sistema};
    \node[item] at (0.35,-7.0)  {$+$ Otros datos};
    \node[item] at (0.35,-7.5)  {Todas las tablas};
    \node[item, font=\tiny] at (0.2,-8.25) {$^{*}$ solo super-administrador};

    \draw[rounded corners] (4.5,0) rectangle (11.4,-8.6);
    \node[anchor=west, font=\small\bfseries] at (4.8,-0.45) {Entradas};
    \node[anchor=west, text width=6.3cm, font=\tiny] at (4.8,-1.0)
      {Resumen: «Artículos del blog. Solo los publicados se muestran en la web.»};
    \draw (4.8,-1.5) rectangle (11.1,-2.0);
    \node[anchor=west, font=\tiny] at (4.9,-1.75) {Buscar\ldots \quad Filtros \quad Columnas \quad Vistas};
    \foreach \y in {-2.4,-2.9,-3.4,-3.9,-4.4} { \draw[gray] (4.8,\y) -- (11.1,\y); }
    \node[anchor=west, font=\tiny] at (4.9,-2.65) {Título \hspace{1.6cm} Estado \hspace{0.9cm} Publicada};
    \draw (4.8,-5.0) rectangle (11.1,-8.3);
    \node[anchor=west, font=\tiny\bfseries] at (4.9,-5.3) {Editar entrada};
    \node[anchor=west, font=\tiny] at (4.9,-5.8) {Título \textcircled{\scalebox{0.7}{?}}};
    \draw (4.9,-6.05) rectangle (10.9,-6.4);
    \node[anchor=west, font=\tiny] at (4.9,-6.8) {Estado \textcircled{\scalebox{0.7}{?}}};
    \draw (4.9,-7.05) rectangle (10.9,-7.4);
    \node[anchor=east, font=\tiny] at (10.9,-7.9) {[Guardar]};
  \end{tikzpicture}
  \caption{Esquema de la consola: navegación por bloques y áreas (izquierda)
    y una pantalla del explorador con su resumen, el listado y un formulario
    con ayuda contextual «?» (derecha).}
  \label{fig:dis-consola}
  % Texto alternativo: a la izquierda, la barra lateral con el bloque Gestión
  % (Inicio, Gestión de cuentas solo para el super-administrador, Usuarios,
  % Archivos, Paneles, Auditoría y Ajustes) y el bloque Datos con las áreas
  % plegables Blog (desplegada: Entradas, Categorías, Etiquetas), Cuentas,
  % Facturación, Sistema y Otros datos, y Todas las tablas. A la derecha, la
  % pantalla de Entradas con un resumen bajo el título, la barra de búsqueda
  % y filtros, el listado y un formulario de edición con un icono de ayuda
  % junto a cada campo.
\end{figure}
\todo{Sustituir o complementar la figura~\ref{fig:dis-consola} con capturas
reales de la consola (pueden generarse con el script de capturas de
\texttt{apps/e2e/scripts}).}

\subsection{Ayuda contextual y resúmenes}
\label{sec:dis-ayudas}

Para que la consola la pueda usar personal no técnico, cada pantalla muestra
bajo el título un \textbf{resumen} de una o dos frases que explica para qué
sirve; en el explorador, ese resumen es la descripción de la tabla guardada en
sus metadatos. Además, cada campo de los formularios de la consola lleva una
\textbf{ayuda contextual «?»}, un componente reutilizable de la biblioteca de
interfaz que muestra al pulsarlo qué significa el campo y qué efecto tiene.
Ambos textos pasan por el sistema de traducción.

\subsection{Internacionalización}
\label{sec:dis-i18n}

La interfaz solo ofrece español (ADR-021): PymeKit se dirige a pymes
españolas y mantener dos idiomas duplicaría el coste de cada texto. Sin
embargo, se conserva toda la infraestructura de traducción: ningún texto
visible se escribe en el código ---tampoco atributos como \texttt{aria-label}
o \texttt{placeholder}---, sino que se referencia por clave a los mensajes del
idioma (\texttt{packages/i18n/src/messages/es}); la lista de idiomas está en
la configuración y una prueba unitaria comprueba la coherencia de las claves.
Añadir un idioma consiste en crear su carpeta de mensajes y registrarlo,
y se documenta como \emph{punto de extensión}. Los identificadores, las
tablas y las columnas siguen en inglés, como es habitual en el ecosistema
(ADR-003); sus nombres visibles en español son datos del CMS.

% ---------------------------------------------------------------------
\section{Decisiones de diseño principales}
\label{sec:dis-decisiones}

Cada decisión relevante se registró en el momento de tomarla con su contexto,
alternativas y consecuencias. La tabla~\ref{tab:dis-adr} las resume; las
decisiones no se borran, sino que se marcan como sustituidas cuando dejan de
valer.

\begin{longtable}{@{}l>{\raggedright\arraybackslash}p{4.4cm}>{\raggedright\arraybackslash}p{6.2cm}l@{}}
  \caption{Resumen de decisiones de diseño (ADR) y alternativas descartadas.}\label{tab:dis-adr}\\
  \toprule
  \textbf{ADR} & \textbf{Decisión} & \textbf{Alternativas descartadas (motivo)} & \textbf{Req.} \\
  \midrule
  \endfirsthead
  \toprule
  \textbf{ADR} & \textbf{Decisión} & \textbf{Alternativas descartadas (motivo)} & \textbf{Req.} \\
  \midrule
  \endhead
  \bottomrule
  \endfoot
  001 & Partir de componentes base con licencia
    & Desarrollo completo desde cero (inviable en plazo); otras bases abiertas (menos completas en multi-tenant y RLS); CMS genéricos con otra base de datos u otro \textit{backend} & RNF-01, 04 \\
  002 & Monorepo único y una sola base de datos & Dos repositorios con la misma base de datos (CI, despliegue y documentación duplicados) & RF-09, RNF-04 \\
  003 & Comentarios y documentación en español; identificadores en inglés & Todo en español (choca con el ecosistema); todo en inglés (dificulta la evaluación) & RF-13, RNF-05 \\
  004 & Un único CMS, de datos, sin CMS de contenidos basado en ficheros & Mantener dos CMS (complejidad sin aportar al objetivo) & RF-09 \\
  005 & Plataforma genérica sin dominio de negocio concreto & Implementar un dominio de ejemplo (desvía el esfuerzo de la base reutilizable) & RNF-01 \\
  006 & Flujo de desarrollo sin comprobaciones remotas de licencia & Mantenerla (dependencia de red y de un servicio externo) & --- \\
  007 & La memoria no nombra los componentes base & --- & --- \\
  008 & Descartar el servidor MCP y los generadores de código ligados a documentación externa & Conservarlos renombrados (dependen de documentación no incluida) & RNF-04 \\
  009 & La web escucha en el puerto 3100 & Puerto por defecto (ocupado en el equipo de desarrollo, B-02) & RNF-07 \\
  010 & Aislar de Internet los servicios locales & Confiar en el cortafuegos del sistema (los contenedores se lo saltan, B-04) & RNF-02 \\
  011 & CMS integrado: API Hono en \texttt{/api/cms}, pantallas como rutas de la web & Aplicaciones separadas (tres servicios, dos inicios de sesión); híbrida (dos enrutadores) & RF-08--11 \\
  012 & Esquema SQL del CMS llamado \texttt{cms} & Conservar el nombre original (excepción permanente en los controles automáticos) & RF-09 \\
  013 & Una sola pila de librerías en toda la web & Mantener las del CMS (dos formas de hacer lo mismo; dos sistemas de i18n) & RNF-05 \\
  014 & Super-administrador como raíz del CMS; RBAC del CMS para el resto del personal & Solo super-administrador, sin RBAC granular (incumple RF-09) & RF-08, 09 \\
  015 & Endurecer la seguridad del CMS & Aceptar el código sin auditar (brechas confirmadas, B-05--B-09) & RNF-02, 03 \\
  016 & Consola para super-administrador y personal del CMS, siempre con \texttt{aal2} & Dejar entrar con el \textit{claim} de acceso sin comprobar el AAL (B-14) & RF-08, RNF-02 \\
  017 & Blog en la base de datos gestionado desde el CMS & Blog fuera (debilita la gestión de contenidos); CMS de ficheros (dos CMS) & RF-01, 09 \\
  018 & Instantánea del autor en la auditoría & Claves ajenas \texttt{RESTRICT} (impide borrar usuarios); borrado lógico (más complejo) & RF-10 \\
  019 & Paneles con cuadrícula propia, sin librería & Librería de cuadrícula (más dependencias; arrastrar sin teclado no es accesible) & RF-11, RNF-08 \\
  020 & Consola por áreas de negocio, nombres en español como datos y sin esquema de demostración & Traducir nombres en el cliente; áreas fijas en un fichero; quitar equipos y pagos & RF-09, RNF-01, 08 \\
  021 & Interfaz solo en español conservando la infraestructura i18n & Mantener dos idiomas (doble mantenimiento); eliminar la i18n (impide traducir) & RF-13, RNF-01 \\
  022 & Una sola imagen Docker para la web y Supabase externa & Contenedores separados para la API del CMS; Supabase autoalojada en el mismo \textit{compose} & RNF-04, 07 \\
  023 & Comentarios didácticos en español limitados al núcleo de la arquitectura & Reescribir todos los comentarios (coste desproporcionado); no reescribir ninguno & RNF-05 \\
  024 & Rol intermedio «Encargado» en los equipos & Mantener solo propietario y miembro; dar al miembro más permisos & RF-06, RNF-01 \\
  025 & Planes de ejemplo para pymes en euros con ids de precio por variables & Ids fijos en el fichero; planes sin precios & RF-07, RNF-01 \\
  026 & Plantilla de roles del CMS en una migración de datos idempotente & Solo en el \textit{seed} (no llega a producción); roles de sistema no editables & RF-09, RNF-01 \\
\end{longtable}

% ---------------------------------------------------------------------
\section{Configuración reutilizable}
\label{sec:dis-configuracion}

La reutilización por configuración (RNF-01) se apoya en tres mecanismos, de
menor a mayor dinamismo.

\paragraph{Ficheros de configuración tipados.} Todo lo que cambia de un
producto a otro se concentra en \texttt{apps/web/src/config/*.config.ts}
(tabla~\ref{tab:dis-config}). Cada fichero valida su contenido con un esquema
Zod al arrancar, de modo que una configuración incorrecta falla pronto y con
un mensaje claro, en lugar de en tiempo de ejecución.

\begin{table}[htbp]
  \centering
  \footnotesize
  \caption{Ficheros de configuración de la aplicación.}
  \label{tab:dis-config}
  \begin{tabularx}{\textwidth}{@{}>{\raggedright\arraybackslash}p{4.4cm}X@{}}
    \toprule
    \textbf{Fichero} & \textbf{Qué configura} \\
    \midrule
    \texttt{app.config.ts} & Nombre del producto, título, descripción, URL y tema por defecto \\
    \texttt{auth.config.ts} & Métodos de acceso (contraseña, enlace mágico, OAuth) y captcha \\
    \texttt{account-mode.config.ts} & Modo de cuentas: personal, equipos o híbrido \\
    \texttt{feature-flags.config.ts} & Interruptores de funcionalidad derivados del modo y de las variables de entorno \\
    \texttt{billing.config.ts} & Proveedor de pago, productos, planes y precios \\
    \texttt{paths.config.ts} & Rutas de la aplicación \\
    \texttt{navigation.config.tsx}, \texttt{marketing-navigation.config.tsx} & Menús del área privada y de la web pública \\
    \texttt{tfg.config.ts} & Bloque académico de la página de inicio \\
    \bottomrule
  \end{tabularx}
\end{table}

\paragraph{Variables de entorno.} Los valores que dependen de la instalación
se leen de variables de entorno. Las públicas llevan el prefijo
\texttt{VITE\_} y los secretos (claves de la base de datos y de la pasarela)
nunca. Los interruptores de los módulos opcionales son:

\begin{itemize}
  \item \textbf{Equipos}: \texttt{VITE\_ACCOUNT\_MODE}
        (\texttt{personal-only}, \texttt{organizations-only} o
        \texttt{hybrid}) decide qué superficies de cuenta existen, y
        \texttt{VITE\_ENABLE\_TEAM\_ACCOUNTS\_CREATION} y
        \texttt{\ldots\_DELETION} las afinan. El modo es la única fuente de
        verdad: un interruptor granular nunca puede reactivar una superficie
        que el modo desactiva.
  \item \textbf{Facturación}: \texttt{VITE\_BILLING\_PROVIDER} y
        \texttt{VITE\_ENABLE\_PERSONAL\_ACCOUNT\_BILLING} /
        \texttt{VITE\_ENABLE\_TEAM\_ACCOUNTS\_BILLING}; la base de datos
        mantiene el mismo interruptor en su tabla de configuración para que
        las políticas RLS también lo respeten.
  \item \textbf{Bloque TFG}: \texttt{VITE\_SHOW\_TFG\_INFO=false} deja la web
        limpia de cualquier referencia académica.
  \item \textbf{Otros}: notificaciones (y su modo en tiempo real), selector
        de tema, baja de cuentas y métodos de autenticación.
\end{itemize}

\paragraph{Configuración como datos.} Lo que la pyme debe poder cambiar sin
desplegar se guarda en la base de datos y se edita desde la consola: el
nombre visible, la descripción y el área de cada tabla, el formato de sus
columnas y el diseño de sus fichas (metadatos del CMS); los roles y permisos
del personal; y las preferencias globales, como la obligatoriedad del MFA.

Finalmente, el monorepo incluye generadores para crear paquetes nuevos con la
estructura y las dependencias correctas, de modo que la lógica propia de un
producto se añade como un paquete más sin modificar el núcleo. El
generador \texttt{package} (\texttt{pnpm turbo gen package}) pregunta si se
trata de una funcionalidad de negocio (\texttt{packages/features}) o de un
paquete de infraestructura y crea su \texttt{package.json},
\texttt{tsconfig.json} y punto de entrada, tomando las versiones del catálogo
del monorepo. Los generadores heredados para otras plataformas de despliegue
se retiraron porque producían configuración de otro \textit{framework}. Toda
esta configuración ---ficheros, variables de entorno, interruptores de
módulos, identidad visual, roles, planes y tablas propias en el CMS--- está
documentada en la guía \texttt{docs/CONFIGURACION.md} del repositorio.
